\chapter{Conclusions}

The Intelligent Compliance Management System (ICMS) demonstrates that integrating Security Information and Event Management (SIEM) telemetry into a Governance, Risk, and Compliance (GRC) engine automates the auditing lifecycle. The platform connects regulatory frameworks to endpoint security through a mapping algorithm. The implementation shows that translating Wazuh Security Configuration Assessment (SCA) alerts into ISO 27001 controls allows an organization to maintain a continuous compliance posture.

The implementation results in several engineering conclusions. A hybrid database schema, using PostgreSQL's \texttt{JSONField} for Wazuh logs alongside relational tables for user roles and frameworks, provides a mechanism for ingesting SIEM data while maintaining ACID compliance. The evaluation phase confirmed that offloading I/O-bound operations, such as API polling and PDF generation, to a Celery worker queue prevents API gateway timeouts. The deployment also shows that stateless JSON Web Tokens (JWT) combined with a React Context provider enforces Role-Based Access Control (RBAC) without server-side session management overhead.

Future development could expand this project into an enterprise product. Technical enhancements include transitioning SIEM integration from REST API polling to a WebSocket or Webhook-driven architecture to reduce latency and enable real-time compliance recalculations. The containerized topology could also be upgraded from Docker Compose to Kubernetes to support horizontal scaling. The long-term objective for the ICMS platform is to refine the orchestration engine for release as an open-source solution. Open-sourcing the codebase would support community development and allow for additional integration modules for SIEM providers like Splunk or Elastic Security, while expanding the library of supported regulatory frameworks.
